\subsection*{Partie A}

Une joueuse de basket constate qu'elle réussit 60\,\% de ses lancers francs au cours d'une séance d'entraînement.

En se basant sur ces statistiques, on modélise la situation comme une succession
d'épreuves aléatoires indépendantes.

\medskip

\begin{enumerate}
	\item Cette joueuse effectue cinq lancers consécutifs.
	\begin{enumerate}
		\item Soit $X$ la variable aléatoire donnant le nombre de lancers réussis.

		Justifier que $X$ suit une loi binomiale dont on précisera les paramètres.
		\item Donner une valeur approchée à $10^{-2}$ près de $P(X \geqslant 4)$.

		Interpréter ce résultat dans le contexte de l'exercice.
	\end{enumerate}
	\item On suppose maintenant que cette sportive doit effectuer $n$ lancers, où $n$ est un entier naturel non nul.

	Soit $Y$ la variable aléatoire donnant le nombre de lancers réussis.

	On admet que la variable aléatoire $Y$ suit la loi binomiale de paramètres $n$ et $p = 0,6$.
	\begin{enumerate}
		\item Montrer que $P(Y \geqslant 1)=1 - 0,4^n$.
		\item En déduire le nombre $n$ de lancers à partir duquel la probabilité d'en
		réussir au moins un dépasse \np{0,9999}.
		\item Recopier et compléter le script de la fonction Python suivante afin qu'elle
		permette de retrouver le résultat de la question \textbf{2. b.}
	\end{enumerate}
\end{enumerate}

\begin{lstlisting}[style=python,xleftmargin=0.3\linewidth,xrightmargin=0.3\linewidth,numbers=none]
def basket ():
	p = 0.6
	n = 1
	while ..... :
		n = n + 1
		p = .....
	return ...
\end{lstlisting}

\subsection*{Partie B}

Pour un autre joueur, on remarque que :
\begin{itemize}
	\item lorsqu'il a réussi un lancer franc, il réussit le suivant avec une probabilité de $0,7$.
	\item lorsqu'il a raté un lancer franc, il réussit le suivant avec une probabilité de $0,4$.
\end{itemize}

On considère que la probabilité de réussir son premier lancer est de  0,5

On note $G_1$ l'évènement : \og le joueur réussit son premier lancer franc \fg.

On a donc $P(G_1) = 0,5$.

On note $G_2$ l'évènement : \og le joueur réussit son deuxième lancer franc \fg .

\medskip

\begin{enumerate}
	\item Calculer $P(G_2)$.

	On pourra utiliser un arbre de probabilité.
\end{enumerate}

Soit $n$ un entier naturel non nul.

On note $G_n$ l'évènement \og le joueur réussit le $n$-ième lancer franc \fg{} et $p_n$, la
probabilité de $G_n$.

On définit ainsi une suite $(p_n)$ de réels.

On admet que pour tout $n \geqslant 1,\: p_{n+1} = 0,3p_n + 0,4$.
\begin{enumerate}[resume]
	\item
	\begin{enumerate}
		\item En utilisant un raisonnement par récurrence, démontrer que

		\begin{center}pour tout $n \geqslant 1,\quad p_n \leqslant p_{n+1}$ \end{center}

		\item Quel argument permet d'affirmer que la suite $(p_n)$ est majorée ?
		\item En déduire que la suite $(p_n)$ est convergente.
	\end{enumerate}
	\item On admet que la limite de $(p_n)$ quand $n$ tend vers $+\infty$ est $\dfrac47$.

	Interpréter cette valeur dans le contexte de l'exercice.
\end{enumerate}
